% =====================================================================
% tesi/capitoli/A7-appendice-crittografia.tex
% =====================================================================
% copre: docs/40-appendice-matematica.md#7-1-cifrario
% copre: docs/40-appendice-matematica.md#7-2-sicurezza-perfetta
% copre: docs/40-appendice-matematica.md#7-3-cifrario-di-vernam-e-chiave-monouso
% copre: docs/40-appendice-matematica.md#7-4-attacco-a-testo-cifrato-noto-e-la-sovrapposizione
% verificato-al-commit: 319226b
% ---------------------------------------------------------------------

\chapter{Nozioni di crittografia}
\label{app:crittografia}

Il capitolo~\ref{cap:cifratura} descrive la cifratura della terza generazione e il capitolo~\ref{cap:analisi} la misura contro un criterio. Questa appendice definisce il criterio e le nozioni che lo circondano, e in particolare rende esplicito ciò che rende la misura significativa: la sicurezza perfetta non è un giudizio comparativo ma una condizione dimostrabile.

\section{Cifrario}

Un cifrario è una coppia di funzioni, cifratura e decifratura, parametrizzate da una chiave, tali che la decifratura con la chiave corretta restituisca il chiaro di partenza. La sicurezza di un cifrario non è proprietà delle sue funzioni ma della difficoltà di risalire al chiaro in assenza della chiave.

La nozione esiste perché consente di distinguere la riservatezza dall'offuscamento, distinzione essenziale nel dominio di questo lavoro e spesso trascurata. La cifratura della terza generazione è un cifrario nel senso pieno della definizione, ma il suo scopo non era proteggere un segreto da un avversario: era rendere costosa la manomissione compiuta con dispositivi di alterazione della memoria.

L'esempio da svolgere è quella medesima struttura, dove la chiave è l'or esclusivo fra valore di personalità e identificativo dell'allenatore e la cifratura è l'or esclusivo di quella chiave con ciascuna parola dei dati. Entrambi i termini della chiave sono presenti in chiaro nella medesima struttura, dunque chi possiede il dato possiede la chiave: il cifrario non protegge da chi legge il salvataggio, e ciò non è una debolezza dell'implementazione ma la conseguenza del suo scopo.

\section{Sicurezza perfetta}

Un cifrario ha sicurezza perfetta, nel senso di Shannon, quando il cifrato non porta alcuna informazione sul chiaro, cioè quando l'informazione mutua fra i due si annulla. La condizione necessaria che ne discende è che l'entropia della chiave sia almeno pari a quella del messaggio.

La nozione esiste perché fornisce un criterio assoluto e non comparativo: non afferma che un cifrario è più robusto di un altro, ma stabilisce se soddisfi una condizione dimostrabile, e la dimostrazione è un confronto fra entropie e non un'analisi di attacchi noti.

L'esempio da svolgere è la misura compiuta nel capitolo~\ref{cap:analisi}. La chiave possiede \SI{32}{bit}, i dati cifrati ne possiedono $384$, cioè le quarantotto parole delle quattro sottostrutture; la condizione richiederebbe almeno $384$ bit di chiave, dunque il deficit vale
\begin{equation}
384 - 32 = 352 \ \text{bit},
\end{equation}
e la sicurezza perfetta non è raggiunta. Va tenuta distinta da questa grandezza il tasso di chiave, cioè il rapporto $384/32 = 12$, che misura quante volte la chiave viene riusata: le due quantità rispondono a domande diverse, e una stesura precedente della nota quantitativa le confondeva. La correzione è registrata perché l'errore è del genere che si ripete: entrambe le grandezze si ottengono dagli stessi due numeri, e soltanto l'enunciato della domanda distingue quale delle due serva.

\section{Cifrario di Vernam e chiave monouso}

Il cifrario di Vernam combina il chiaro con una chiave della medesima lunghezza mediante or esclusivo. Prende il nome di chiave monouso quando la chiave è casuale, lunga quanto il messaggio e impiegata una sola volta, e sotto quelle tre condizioni raggiunge la sicurezza perfetta della sezione precedente.

La nozione esiste perché è l'unico cifrario la cui sicurezza sia dimostrata anziché congetturata, e perché le sue tre condizioni costituiscono il metro con cui si misura ogni cifrario a chiave scorrevole reale: ciascuna, violata, apre un attacco specifico e noto.

L'esempio da svolgere è la violazione che riguarda questo lavoro. La cifratura della terza generazione soddisfa la forma del cifrario di Vernam ma viola la seconda e la terza condizione, poiché la chiave è dodici volte più corta del messaggio e viene riusata su ciascuna parola. Ne discende l'attacco della sezione seguente, che non richiede di indovinare alcuna chiave.

\section{Attacco a testo cifrato noto, e la sovrapposizione}

Un attacco a testo cifrato noto è condotto da chi possiede il solo cifrato, senza alcun chiaro corrispondente. È il modello di avversario più debole e quindi il più realistico, e un cifrario che ceda contro di esso non offre protezione pratica.

La nozione esiste perché la robustezza di un cifrario si dichiara sempre rispetto a un modello di avversario, e omettere il modello priva l'affermazione di contenuto: un cifrario può resistere a chi possieda il solo cifrato e cedere contro chi disponga di una coppia di chiaro e cifrato.

L'esempio da svolgere è l'attacco per sovrapposizione esposto nel capitolo~\ref{cap:analisi}, di cui qui si riporta la struttura logica. Se due parole $d_1$ e $d_2$ sono cifrate con la medesima chiave $k$, l'avversario che osservi $c_1 = d_1 \oplus k$ e $c_2 = d_2 \oplus k$ calcola
\begin{equation}
c_1 \oplus c_2 = (d_1 \oplus k) \oplus (d_2 \oplus k) = d_1 \oplus d_2,
\end{equation}
ottenendo una relazione fra i chiari senza conoscere la chiave. Poiché molte parole delle sottostrutture sono nulle o assumono valori prevedibili, quella relazione rivela direttamente il contenuto delle parole ignote. Il difetto non risiede nell'or esclusivo, che è la forma corretta del cifrario di Vernam, ma nel riuso della chiave, cioè nella seconda condizione violata della sezione precedente: è la distinzione che separa un algoritmo mal scelto da un algoritmo bene scelto e mal impiegato.
